% =====================================================================
%  chapters/minggu-14.tex
%  Minggu 14 — Workshop Proyek Akhir: Riset UI/UX & Perencanaan Desain
% =====================================================================

\chapter{Workshop Proyek Akhir --- Riset UI/UX \& Perencanaan Desain}
\label{chap:minggu-14}

% ---------------------------------------------------------------------
% 1. CAPAIAN PEMBELAJARAN
% ---------------------------------------------------------------------
\begin{capaian}
Setelah menyelesaikan bab ini, mahasiswa diharapkan mampu:
\begin{itemize}
  \item melakukan \emph{heuristic evaluation} terhadap halaman website \emph{e-governance} menggunakan 10 \emph{Usability Heuristics} (Nielsen Norman Group);
  \item mendokumentasikan \emph{pain point} dan hasil \emph{benchmarking} secara terstruktur;
  \item membuat \emph{wireframe} dan \emph{mockup} desain ulang di Figma berdasarkan temuan riset;
  \item menyusun rasional desain yang menghubungkan temuan riset dengan keputusan desain.
\end{itemize}
\end{capaian}

% ---------------------------------------------------------------------
% 2. TUJUAN PRAKTIK
% ---------------------------------------------------------------------
\begin{tujuanpraktik}
Pada sesi praktik ini, mahasiswa akan:
\begin{itemize}
  \item menyiapkan struktur \emph{repository} dan \emph{branch} untuk proyek akhir;
  \item mengidentifikasi halaman website target dan mendokumentasikannya;
  \item melakukan \emph{heuristic evaluation} dan \emph{benchmarking};
  \item membuat \emph{wireframe} dan \emph{mockup} di Figma;
  \item menulis rasional desain.
\end{itemize}
\end{tujuanpraktik}

% ---------------------------------------------------------------------
% 3. PERSIAPAN KELOMPOK
% ---------------------------------------------------------------------
\section{Persiapan Kelompok}
\label{sec:persiapan-kelompok}

Setiap kelompok terdiri atas 3 orang. Berikut peralatan yang dibutuhkan:

\begin{itemize}
  \item \textbf{Browser} (Chrome/Firefox) --- akses website target dan \emph{heuristic checklist}.
  \item \textbf{Figma} (akun gratis) --- \emph{wireframing} dan \emph{mockup}.
  \item \textbf{VS Code} --- persiapan untuk minggu 15.
  \item \textbf{GitHub} --- \emph{repository} proyek akhir.
  \item \textbf{Google Docs / Markdown editor} --- dokumentasi hasil riset.
\end{itemize}

\subsection{Inisialisasi Repository}
\label{subsec:inisialisasi-repo}

\begin{langkah}
  \item Buka \emph{repository} GitHub kelompok (yang sudah dibuat dari minggu 8--9 atau buat baru).
  \item Buat \emph{branch} baru untuk proyek akhir:
\begin{lstlisting}[language=bash, caption={Membuat branch final-project}]
git checkout -b final-project
\end{lstlisting}
  \item Buat struktur folder di dalam \emph{repo}:
\begin{lstlisting}[caption={Struktur folder proyek akhir}]
final-project/
  research/
    heuristic-evaluation.md
    benchmarking.md
  design/
  src/
    index.html
    css/
    js/
    assets/
  README.md
\end{lstlisting}
  \item \emph{Push branch} ke GitHub:
\begin{lstlisting}[language=bash, caption={Push awal}]
git add .
git commit -m "feat: inisialisasi struktur folder proyek akhir"
git push -u origin final-project
\end{lstlisting}
\end{langkah}

\begin{tangkapanlayar}
  {assets/images/minggu-14/struktur-folder-repo.png}
  {Struktur folder repo di VS Code Explorer.}
  {Buka \texttt{VS Code} $\rightarrow$ \texttt{Explorer} $\rightarrow$ periksa bahwa seluruh folder dan file sudah sesuai.}
\end{tangkapanlayar}

% ---------------------------------------------------------------------
% 4. IDENTIFIKASI WEBSITE TARGET
% ---------------------------------------------------------------------
\section{Identifikasi Website Target}
\label{sec:website-target}

\begin{langkah}
  \item Pilih satu halaman dari website \emph{e-governance} / layanan publik, contoh:
  \begin{itemize}
    \item Portal layanan perizinan daerah.
    \item Halaman pendaftaran online instansi pemerintah.
    \item Portal informasi desa/kelurahan.
    \item Halaman pengaduan masyarakat \emph{online}.
  \end{itemize}
  \item Dokumentasikan informasi umum di \texttt{research/heuristic-evaluation.md} --- nama website, URL, halaman yang diriset, tanggal, dan anggota kelompok beserta peran masing-masing.
  \item Ambil \emph{screenshot} halaman asli dari berbagai ukuran layar (desktop dan mobile) menggunakan \emph{Developer Tools} (F12) $\rightarrow$ \emph{toggle device toolbar}.
\end{langkah}

\begin{tangkapanlayar}
  {assets/images/minggu-14/halaman-target-desktop-mobile.png}
  {Tampilan halaman website target --- versi desktop dan mobile.}
  {Buka halaman target di browser, ambil \emph{screenshot} versi desktop, lalu gunakan \emph{Device Toolbar} di \emph{DevTools} untuk versi mobile.}
\end{tangkapanlayar}

\begin{tip}
Jika sulit menentukan website target, pilih website pemerintah daerah (portal layanan disdukcapil, BPJS, atau website desa). Hindari website yang sudah sangat modern.
\end{tip}

% ---------------------------------------------------------------------
% 5. HEURISTIC EVALUATION
% ---------------------------------------------------------------------
\section{Heuristic Evaluation}
\label{sec:heuristic-evaluation}

Gunakan 10 \emph{Usability Heuristics} dari Nielsen Norman Group berikut sebagai panduan evaluasi.

\begin{table}[htbp]
\centering
\caption{10 Usability Heuristics --- Nielsen Norman Group}
\label{tab:heuristik-nng}
\small
\begin{tabular}{@{}clp{7cm}@{}}
\toprule
\textbf{\#} & \textbf{Heuristik} & \textbf{Pertanyaan Panduan} \\
\midrule
1 & Visibility of System Status & Apakah pengguna tahu apa yang sedang terjadi di halaman? \\
2 & Match Between System and Real World & Apakah bahasa dan istilah sesuai dunia nyata pengguna? \\
3 & User Control and Freedom & Apakah pengguna bisa membatalkan aksi atau navigasi dengan mudah? \\
4 & Consistency and Standards & Apakah tampilan dan interaksi konsisten di seluruh halaman? \\
5 & Error Prevention & Apakah ada perlindungan untuk mencegah kesalahan pengguna? \\
6 & Recognition Rather Than Recall & Apakah pengguna perlu mengingat informasi dari satu bagian ke bagian lain? \\
7 & Flexibility and Efficiency of Use & Apakah ada cara cepat untuk pengguna yang sudah terbiasa? \\
8 & Aesthetic and Minimalist Design & Apakah antarmuka bebas dari elemen yang tidak perlu? \\
9 & Help Users Recognize, Diagnose, and Recover from Errors & Apakah pesan \emph{error} jelas dan membantu? \\
10 & Help and Documentation & Apakah ada panduan atau bantuan jika pengguna kebingungan? \\
\bottomrule
\end{tabular}
\end{table}

\subsection{Langkah Evaluasi}
\label{subsec:langkah-evaluasi}

Untuk setiap heuristik, diskusikan secara berkelompok dan catat di \texttt{research/heuristic-evaluation.md}. Tuliskan: skor pelanggaran (0--4, di mana 0 = tidak ada masalah dan 4 = masalah kritis), temuan, \emph{screenshot} bukti, dampak ke pengguna, serta saran perbaikan yang akan diterapkan di desain ulang.

\begin{langkah}
  \item Buka \emph{heuristic checklist} di atas satu per satu.
  \item Untuk setiap heuristik, identifikasi apakah ada pelanggaran pada halaman target.
  \item Catat temuan beserta \emph{screenshot} bukti.
  \item Uraikan dampak temuan terhadap pengalaman pengguna.
  \item Hubungkan setiap temuan dengan rekomendasi desain yang akan diterapkan.
\end{langkah}

\begin{tip}
Jangan khawatir soal skor yang ``benar'' --- yang penting adalah kemampuan menjelaskan mengapa suatu masalah muncul. Satu heuristik bisa memiliki 0 atau lebih temuan. Jika tidak ada masalah, tulis ``Tidak ditemukan pelanggaran signifikan.''
\end{tip}

\subsection{Ringkasan Temuan}
\label{subsec:ringkasan-temuan}

Buat tabel ringkasan di bagian bawah file \texttt{heuristic-evaluation.md} berisi nomor heuristik, skor, dan jumlah temuan. Identifikasi 3 \emph{pain point} terbesar dari seluruh evaluasi.

\begin{tangkapanlayar}
  {assets/images/minggu-14/heuristic-evaluation-dokumentasi.png}
  {Dokumentasi heuristic evaluation di Markdown --- tampilan di VS Code atau GitHub.}
  {Buka file \texttt{research/heuristic-evaluation.md} di VS Code, periksa bahwa seluruh heuristik sudah dievaluasi dan tabel ringkasan sudah dilengkapi.}
\end{tangkapanlayar}

% ---------------------------------------------------------------------
% 6. BENCHMARKING
% ---------------------------------------------------------------------
\section{Benchmarking Singkat}
\label{sec:benchmarking}

Bandingkan halaman target dengan minimal 2--3 website sejenis yang memiliki fitur serupa.

\begin{langkah}
  \item Pilih minimal 2 website untuk dibandingkan (misalnya website pemerintah daerah lain atau layanan publik modern).
  \item Buat file \texttt{research/benchmarking.md} yang memuat:
  \begin{itemize}
    \item Informasi setiap website: URL, halaman sejenis, kelebihan, kekurangan, dan \emph{screenshot}.
    \item Tabel perbandingan fitur antara website target, website pembanding, dan rencana desain ulang.
    \item Daftar pelajaran dari \emph{benchmarking} (fitur yang ingin diadopsi, yang perlu dihindari, dan peluang diferensiasi).
  \end{itemize}
\end{langkah}

\begin{tangkapanlayar}
  {assets/images/minggu-14/benchmarking-perbandingan.png}
  {Perbandingan visual antara website target dan website pembanding.}
  {Buka file \texttt{research/benchmarking.md} dan pastikan tabel perbandingan fitur sudah terisi lengkap.}
\end{tangkapanlayar}

% ---------------------------------------------------------------------
% 7. PERENCANAAN DESAIN DI FIGMA
% ---------------------------------------------------------------------
\section{Perencanaan Desain di Figma}
\label{sec:desain-figma}

\subsection{Membuat Frame Wireframe}
\label{subsec:wireframe}

\begin{langkah}
  \item Buka Figma (\emph{desktop app} atau browser), klik \textbf{``New design file''}.
  \item Buat \emph{frame} --- tekan \texttt{F} lalu pilih ukuran: Desktop 1440$\times$900, Mobile 375$\times$812. Beri nama ``Wireframe --- Desktop'' dan ``Wireframe --- Mobile''.
  \item Sketsa tata letak \textbf{tanpa warna}, hanya menggunakan:
  \begin{itemize}
    \item Kotak untuk area konten.
    \item Garis atau garis putus untuk teks \emph{placeholder}.
    \item Ikon sederhana untuk elemen navigasi.
    \item Anotasi singkat di sebelah setiap elemen.
  \end{itemize}
\end{langkah}

\begin{tip}
Gunakan hanya warna abu-abu, hitam, dan putih pada wireframe. Jangan tambahkan warna, gambar, atau font khusus. Fokus pada struktur dan alur pengguna, bukan visual.
\end{tip}

\begin{tangkapanlayar}
  {assets/images/minggu-14/wireframe-desktop-mobile.png}
  {Wireframe desktop dan mobile di Figma --- tampilan keseluruhan.}
  {Buka file Figma, periksa bahwa \emph{frame} wireframe sudah dibuat untuk versi desktop dan mobile dengan anotasi yang jelas.}
\end{tangkapanlayar}

\subsection{Membuat Mockup}
\label{subsec:mockup}

\begin{langkah}
  \item Duplikat \emph{frame} wireframe $\rightarrow$ rename menjadi ``Mockup --- Desktop'' dan ``Mockup --- Mobile''.
  \item Terapkan visual pada \emph{mockup}:
  \begin{itemize}
    \item Warna merek (maksimal 3--4 warna).
    \item Font yang sudah dipilih (\emph{Google Fonts} atau \emph{font system}).
    \item Gambar/ikon asli.
    \item Spasi dan tipografi yang konsisten.
  \end{itemize}
  \item Pastikan konsistensi antar halaman --- header/footer identik, tipografi konsisten, spasi antar elemen seragam.
\end{langkah}

\begin{tangkapanlayar}
  {assets/images/minggu-14/mockup-desktop-overview.png}
  {Mockup desktop di Figma --- tampilan keseluruhan.}
  {Periksa bahwa \emph{mockup} sudah menerapkan warna merek, tipografi, dan spasi secara konsisten.}
\end{tangkapanlayar}

\begin{tangkapanlayar}
  {assets/images/minggu-14/mockup-mobile-overview.png}
  {Mockup mobile di Figma --- tampilan keseluruhan.}
  {Periksa bahwa versi mobile memiliki tata letak yang responsif dan konsisten dengan versi desktop.}
\end{tangkapanlayar}

\subsection{Menulis Rasional Desain}
\label{subsec:rasional-desain}

\begin{langkah}
  \item Buat file \texttt{research/rasional-desain.md} atau tulis langsung di \texttt{README.md}.
  \item Isi tabel yang menghubungkan temuan riset dengan keputusan desain: nomor, temuan riset, keputusan desain, dan penjelasan mengapa keputusan tersebut diambil.
  \item Daftar prinsip desain yang diterapkan (misalnya \emph{Consistency}, \emph{Minimalism}, \emph{Error Prevention}).
\end{langkah}

\begin{catatan}
Kaitan antara temuan riset dan keputusan desain merupakan komponen kunci yang akan dinilai pada presentasi minggu 16. Pastikan setiap keputusan desain dapat ditelusuri kembali ke minimal satu temuan riset.
\end{catatan}

% ---------------------------------------------------------------------
% 8. COMMIT \& PUSH
% ---------------------------------------------------------------------
\section{Commit dan Push}
\label{sec:commit-push}

\begin{langkah}
  \item Kembali ke VS Code, buka terminal.
  \item Pastikan semua file sudah tersimpan di folder yang benar.
  \item \emph{Commit} dan \emph{push}:
\begin{lstlisting}[language=bash, caption={Commit hasil riset}]
git add .
git commit -m "feat: hasil riset heuristic evaluation, benchmarking, dan wireframe/mockup"
git push
\end{lstlisting}
  \item Buka \emph{pull request} di GitHub (opsional, tetapi disarankan) dengan deskripsi singkat mengenai isi \emph{commit}.
\end{langkah}

\begin{tangkapanlayar}
  {assets/images/minggu-14/commit-history-github.png}
  {Commit history di GitHub --- terlihat branch final-project.}
  {Buka \emph{repository} di GitHub, navigasi ke tab ``Commits'' untuk memastikan \emph{commit} berhasil.}
\end{tangkapanlayar}

% ---------------------------------------------------------------------
% 9. OUTPUT YANG DIHARAPKAN
% ---------------------------------------------------------------------
\section{Output yang Diharapkan}
\label{sec:output-minggu-14}

Setelah sesi ini, setiap kelompok harus memiliki:
\begin{itemize}
  \item Dokumen \emph{heuristic evaluation} (\texttt{research/heuristic-evaluation.md}).
  \item Dokumen \emph{benchmarking} (\texttt{research/benchmarking.md}).
  \item \emph{Wireframe} (desktop \& mobile) di Figma.
  \item \emph{Mockup} (desktop \& mobile) di Figma.
  \item Rasional desain (\texttt{research/rasional-desain.md} atau \texttt{README.md}).
  \item \emph{Screenshot} halaman asli di folder \texttt{assets/} atau \texttt{research/}.
  \item \emph{Commit} di GitHub pada branch \texttt{final-project}.
\end{itemize}

% ---------------------------------------------------------------------
% 10. TROUBLESHOOTING
% ---------------------------------------------------------------------
\section{Troubleshooting}
\label{sec:troubleshooting-minggu-14}

\begin{table}[htbp]
\centering
\caption{Masalah Umum dan Solusi --- Minggu 14}
\label{tab:troubleshooting-14}
\small
\begin{tabular}{@{}p{5.5cm}p{8cm}@{}}
\toprule
\textbf{Masalah} & \textbf{Solusi} \\
\midrule
Website target tidak bisa diakses & Gunakan Wayback Machine (\texttt{web.archive.org}) untuk versi terakhir yang bisa diakses, atau pilih website lain. \\
Tidak yakin bagaimana mengisi \emph{heuristic evaluation} & Lihat contoh di Nielsen Norman Group --- fokus pada kemampuan menjelaskan, bukan skor sempurna. \\
Figma tidak bisa dibuka di browser & Gunakan \emph{Figma desktop app} (download dari \texttt{figma.com}), atau pastikan browser versi terbaru. \\
Anggota kelompok di lokasi berbeda & Semua file riset ada di GitHub --- akses dari mana saja. Desain Figma bisa di-\emph{share} via tautan. \\
Tidak bisa membuat \emph{branch} di Git & Pastikan sudah \emph{clone} \emph{repository} ke lokal. Jika masih gagal, cek \texttt{git status}. \\
Pekerjaan terasa terlalu banyak untuk 1 pertemuan & Fokus pada \emph{heuristic evaluation} dan \emph{wireframe} terlebih dahulu. \emph{Mockup} bisa dilanjutkan di minggu 15. \\
\bottomrule
\end{tabular}
\end{table}

% ---------------------------------------------------------------------
% 11. TAUTAN REFERENSI
% ---------------------------------------------------------------------
\section{Tautan Referensi}
\label{sec:referensi-minggu-14}

\begin{itemize}
  \item 10 Usability Heuristics (NNGroup): \href{https://www.nngroup.com/articles/ten-usability-heuristics/}{nngroup.com/articles/ten-usability-heuristics}
  \item Figma Getting Started: \href{https://help.figma.com/hc/en-us/articles/360040319993}{help.figma.com/hc/en-us/articles/360040319993}
  \item Heuristic Evaluation Template (NNGroup): \href{https://www.nngroup.com/articles/how-to-conduct-a-heuristic-evaluation/}{nngroup.com/articles/how-to-conduct-a-heuristic-evaluation}
  \item Wayback Machine: \href{https://web.archive.org}{web.archive.org}
  \item Contoh Website E-Governance: \url{https://www.go.id}, \url{https://jakartasatu.jakarta.go.id}
\end{itemize}

% ---------------------------------------------------------------------
% 12. PERSIAPAN UNTUK MINGGU 15
% ---------------------------------------------------------------------
\section{Persiapan untuk Minggu 15}
\label{sec:persiapan-minggu-15}

Sebelum pertemuan minggu depan, pastikan:
\begin{enumerate}
  \item Semua dokumen riset sudah \emph{ter-commit} dan \emph{ter-push} ke GitHub.
  \item Desain Figma sudah selesai (\emph{wireframe} + \emph{mockup}).
  \item Anggota kelompok sudah berbagi peran untuk implementasi:
  \begin{itemize}
    \item Anggota A: header, navigasi, \emph{hero section}.
    \item Anggota B: bagian tengah (konten utama, formulir, informasi).
    \item Anggota C: footer, bagian bawah, optimasi responsif.
  \end{itemize}
  \item File \texttt{index.html} sudah mulai disiapkan (kosong atau \emph{skeleton} HTML sudah ada di branch \texttt{final-project}).
\end{enumerate}

% ---------------------------------------------------------------------
% 13. LATIHAN
% ---------------------------------------------------------------------
\section{Latihan Mandiri}
\label{sec:latihan-minggu-14}

\begin{latihan}[Eksplorasi Heuristik]
\begin{enumerate}[label=\alph*., leftmargin=2em, itemsep=2pt]
  \item Pilih salah satu heuristik dari tabel, lalu buka 3 website berbeda. Catat minimal 1 temuan per website yang berkaitan dengan heuristik tersebut.
  \item Bandingkan kemudahan penggunaan antara kedua website --- tuliskan alasan Anda secara singkat.
\end{enumerate}
\end{latihan}
